\DocumentMetadata{
  pdfversion=2.0,
  pdfstandard=ua-2,
  lang=en-US,
  tagging=on,
  tagging-setup={math/setup=mathml-SE,tabsorder=structure}
}
\documentclass[11pt,letterpaper]{article}
\usepackage[margin=0.68in,includefoot]{geometry}
\usepackage{newcomputermodern}
\usepackage{amsmath,amscd,mathtools}
\usepackage{tikz,adjustbox}
\providecommand{\square}{\mdlgwhtsquare}
\usepackage[hidelinks]{hyperref}
\hypersetup{
  pdftitle={Numerical Linear Algebra First-Year Exam, January 2022},
  pdfauthor={University of Florida Department of Mathematics},
  pdfsubject={Graduate mathematics examination},
  pdfdisplaydoctitle=true,
  bookmarksopen=true
}
\setcounter{secnumdepth}{0}
\setlength{\parindent}{0pt}
\setlength{\parskip}{0.28em}
\setlength{\emergencystretch}{3em}
\setlength{\leftmargini}{2.15em}
\setlength{\leftmarginii}{2.65em}
\setlength{\leftmarginiii}{3em}
\pagestyle{empty}
\raggedbottom
\allowdisplaybreaks
\NewTaggingSocketPlug{block/list/label}{examproblem}{%
  \tagstructbegin{tag=Lbl}\tagstructend%
  \tagstructbegin{tag=\LBody}%
  \tagstructbegin{tag=\ProblemHeadingTag,title-o={\ProblemHeadingText}}%
  \tagmcbegin{tag=\ProblemHeadingTag,actualtext={\ProblemHeadingText}}%
  #2%
  \tagmcend%
  \tagstructend%
}
\newenvironment{examproblems}[2]{%
  \def\ProblemHeadingTag{#2}%
  \AssignTaggingSocketPlug{block/list/label}{examproblem}%
  \begin{enumerate}%
  \setcounter{enumi}{#1}%
  \renewcommand{\labelenumi}{\textbf{\theenumi.}}%
  \setlength{\topsep}{0.35em}%
  \setlength{\partopsep}{0pt}%
  \setlength{\itemsep}{0.48em}%
  \setlength{\parsep}{0.18em}%
}{\end{enumerate}}
\newenvironment{examsubparts}{%
  \AssignTaggingSocketPlug{block/list/label}{default}%
  \begin{enumerate}%
  \setlength{\topsep}{0.18em}%
  \setlength{\partopsep}{0pt}%
  \setlength{\itemsep}{0.16em}%
  \setlength{\parsep}{0.1em}%
}{\end{enumerate}}
\newenvironment{examinstructions}{%
  \par\begingroup\itshape
}{%
  \par\endgroup\vspace{0.18em}
}
\makeatletter
\renewcommand\subsection{\@startsection{subsection}{2}{0pt}%
  {-1.25ex plus -.25ex minus -.1ex}%
  {0.55ex plus .1ex}%
  {\normalfont\large\bfseries}}
\renewcommand{\@maketitle}{%
  \newpage\null\vskip -1em%
  {\footnotesize\bfseries
    \href{https://www.ufl.edu/}{University of Florida}, \href{https://math.ufl.edu/}{Department of Mathematics}\par}%
  \vskip -0.5em%
  \tagstructbegin{tag=H1,title={Numerical Linear Algebra First-Year Exam, January 2022}}%
  \tagpdfparaOff%
  \tagmcbegin{tag=H1}%
    {\LARGE\bfseries\href{https://gma.math.ufl.edu/exams/numerical-linear-algebra-first-year/}{Numerical Linear Algebra First-Year Exam}, January 2022\par}%
  \tagmcend%
  \tagpdfparaOn%
  \tagstructend%
  \vskip 0.25em%
  \tagmcbegin{artifact}%
    \hrule height 1pt%
  \tagmcend%
  \vskip 1em%
}
\makeatother
\title{Numerical Linear Algebra First-Year Exam, January 2022}
\author{}
\date{}
\begin{document}
\maketitle
\thispagestyle{empty}
\begin{examinstructions}
Do 4 (four) problems.
\end{examinstructions}
\begin{examproblems}{0}{H2}
\def\ProblemHeadingText{Problem 1.}
\item
Let \(A \in \mathbb{C}^{m\times n}\) and \(B \in \mathbb{C}^{n\times p}\). Prove or provide a counterexample to the following statements. Justify each nontrivial step.
\begin{examsubparts}
\item[\textnormal{(a)}]
The Frobenius norm satisfies \(\lVert AB\rVert_F \leq \lVert A\rVert_F\lVert B\rVert_F\).
\item[\textnormal{(b)}]
\(\lVert A\rVert_2 \leq \lVert A\rVert_F\), where \(\lVert A\rVert_F\) is the Frobenius norm of~\(A\).
\end{examsubparts}
\def\ProblemHeadingText{Problem 2.}
\item
Define the matrices~\(A\) and~\(B\) by
\[
A=\begin{pmatrix}1&2&3\\2&4&6\\1&2&3\\2&4&6\end{pmatrix},\qquad B=\begin{pmatrix}1&1&0\\0&1&3\\0&1&3\\2&0&0\end{pmatrix}.
\]
\begin{examsubparts}
\item[\textnormal{(a)}]
Find an economy singular value decomposition of~\(A\).
\item[\textnormal{(b)}]
Find an economy QR decomposition of~\(B\).
\end{examsubparts}
\def\ProblemHeadingText{Problem 3.}
\item
Let \(\lVert\cdot\rVert\) be a subordinate (induced) matrix norm.
\begin{examsubparts}
\item[\textnormal{(a)}]
If~\(E\) is \(n\times n\) with \(\lVert E\rVert<1\), then show \(I+E\) is nonsingular and
\[
\lVert(I+E)^{-1}\rVert\leq \frac{1}{1-\lVert E\rVert}.
\]
\item[\textnormal{(b)}]
If~\(A\) is \(n\times n\) invertible and~\(E\) is \(n\times n\) with \(\lVert A^{-1}\rVert\lVert E\rVert<1\), then show \(A+E\) is nonsingular and
\[
\lVert(A+E)^{-1}\rVert\leq \frac{\lVert A^{-1}\rVert}{1-\lVert A^{-1}\rVert\lVert E\rVert}.
\]
\end{examsubparts}
\def\ProblemHeadingText{Problem 4.}
\item
Let \(P\in\mathbb{C}^{m\times m}\) be a projector. Show that \(\lVert P\rVert_2=1\) if and only if~\(P\) is an orthogonal projector.
\def\ProblemHeadingText{Problem 5.}
\item
Suppose~\(A\) is Hermitian positive definite.
\begin{examsubparts}
\item[\textnormal{(a)}]
Prove that each principal submatrix of~\(A\) is Hermitian positive definite.
\item[\textnormal{(b)}]
Prove that an element of~\(A\) with largest magnitude lies on the diagonal.
\item[\textnormal{(c)}]
Prove that~\(A\) has a Cholesky decomposition.
\end{examsubparts}
\end{examproblems}
\end{document}
